\chapter{Epidemias}


A lo largo de toda nuestra historia como especie, nos hemos enfrentado a problemas de salud, dentro de los cuales los virus y las bacterias han sido parte fundamental de estas situaciones. No importa cuánto logremos desarrollar la medicina: así como nosotros evolucionamos como especie, también lo hace nuestro ambiente, por lo que estamos en una lucha constante por sobrevivir.

Recientemente, el mundo se enfrentó a la pandemia de COVID-19, un evento que impactó al mundo entero y que, a pesar de todo el conocimiento que se tiene sobre el área, vino a demostrar que es necesario seguir desarrollando temas epidemiológicos para poder afrontar este tipo de problemas.

Es importante notar que este evento no ha sido el único del que tenemos registro. Desde la antigüedad, la humanidad se ha enfrentado a diversas epidemias, siendo una de las más antiguas de las que se tiene registro la plaga de Atenas, que data aproximadamente de los años 430 a 426 a.\,C.

Fue gracias a Tucídides, militar e historiador ateniense, que se tiene un registro detallado de lo ocurrido en esta epidemia. Él fue quien recabó distinta información sobre la enfermedad que azotó al pueblo ateniense, registrada en un texto en el que se lee:

\begin{quote}
«Descrito con gran detalle [\dots] incluidos los síntomas y la progresión de la enfermedad. También dio algunas cifras exactas, como que 1050 de 4000 soldados en una expedición murieron a causa de la enfermedad. La enfermedad descrita tan minuciosamente por Tucídides, incluso hasta el hecho de que los perros que comieron los cadáveres también sufrieron, ha sido la fuente de numerosos artículos durante cientos de años con casos presentados (con gran convicción y defendidos vehementemente) para un increíble rango de enfermedades».\cite{murray1989}
\end{quote}

Otro caso, posiblemente el más famoso después del COVID-19, es el de la peste negra, la cual:

\begin{quote}
«se convirtió en una pandemia de tal magnitud que no solo sacudió los cimientos económicos y sociales del Viejo Mundo, sino que cambió el curso de la historia. Tuvo su origen en Asia, se propagó a Europa, y su influencia llegó hasta el norte de África. El contagio llegó por barco, pues en aquellos tiempos empezaron a abrirse nuevas rutas de intercambio, ya que Europa empezaba a tener un avance económico y comercial».\cite{german2021}
\end{quote}

Se estima que la peste negra causó entre 75 y 200 millones de muertes en el mundo, mientras que en Europa murió entre el 30\,\% y el 60\,\% de la población.

Estos y demás ejemplos que podemos seguir tratando nos dan una idea de por qué es importante estudiar este tipo de fenómenos, ya que las implicaciones sociales, culturales, entre otras, son de gran relevancia para el bienestar colectivo.

\section{¿Cómo se comporta una epidemia?}

Ahora que tenemos el contexto sobre la importancia de estudiar las epidemias, podemos plantearnos varias preguntas que nos ayudarán a abordar esta área.

Definamos los siguientes conceptos, que nos darán la pauta para establecer las reglas que seguirá nuestro estudio.

Una epidemia es ``cuando el número de casos de una enfermedad excede lo esperado en una comunidad, región o país, donde se especifican claramente la región geográfica y el periodo temporal'' \cite{pelaez2020}.

Una infección, en términos generales, es ``la invasión de un organismo por parte de otro más pequeño'' \cite{vynnycky2010}. Con esta definición podemos hacer referencia no únicamente a los casos de enfermedades en humanos, sino también a que la infección es una situación común en la naturaleza.

Entonces, lo importante es entender la dinámica de una pandemia. Con lo definido anteriormente, se trata de preguntarnos cómo ocurre la transición del estado de una persona no infectada a un estado de infectada, y esto a gran escala.

Dada la naturaleza de la enfermedad que se estudia, la cantidad de suposiciones que se tomarán en cuenta es considerable, ya que no es posible tratar todos los casos de la misma manera. Los modelos clásicos trabajan a partir de una cantidad fija de personas susceptibles, por lo que casos como nacimientos o muertes que no sean causados por la enfermedad no se contabilizan.

Esto nos lleva a preguntarnos: ¿cuál es la metodología para elegir un modelo apropiado? Para responderla, debemos analizar las características de nuestro caso de estudio e identificar cuáles son cruciales y cuáles podemos dejar de lado. En el caso de las enfermedades, algunas de las preguntas que podemos plantearnos para establecer el modelo son:

\begin{itemize}
    \item ¿Es muy contagiosa esta enfermedad?
    \item ¿La enfermedad tiene una alta mortalidad?
    \item ¿Quiénes son propensos a contraer la enfermedad?
    \item ¿Qué tan rápido se desarrolla la enfermedad en una persona?
    \item Una persona infectada, ¿cuánto tiempo tiene la enfermedad?
\end{itemize}

A partir de las preguntas adecuadas y del objetivo del estudio, se puede elegir un modelo correcto. Dependiendo de los objetivos, podemos optar entre un modelo determinista o uno estocástico. Los modelos deterministas incluyen información sobre los promedios de lo que sucede en la población. Por otro lado, los modelos estocásticos incorporan la aleatoriedad del proceso y nos brindan información numérica sobre las transiciones entre diferentes compartimentos.

La mayoría de los modelos deterministas son comúnmente llamados modelos compartimentales. Estos dividen la población en subgrupos y, mediante sistemas de ecuaciones diferenciales, se aproxima la cantidad de personas que se encuentran en cada subgrupo. Para ello, el tiempo se discretiza en pasos, generalmente de un día, lo que permite obtener la aproximación del tiempo $t+1$ en términos de lo que se tiene en $t$.

\section{El modelo SIRD}

El modelo SIRD es un modelo compartimental en el que se divide una población total $N$ (constante) en cuatro compartimientos:

\begin{itemize}
    \item Susceptibles ($S$)
    \item Infectados ($I$)
    \item Recuperados ($R$)
    \item Fallecidos ($D$)
\end{itemize}

Cabe notar que, a diferencia de modelos como SIR o SEIR, aquí se toman en cuenta las personas fallecidas, ya que la mortalidad registrada durante la pandemia de COVID-19 fue considerable y no es posible despreciarla.

Además, estos compartimientos seguirán las siguientes reglas para el estudio que se presentará más adelante:

\begin{itemize}
    \item El compartimiento de susceptibles incluye a cualquier persona que pueda ser infectada.
    \item El compartimiento de infectados incluye a cualquier persona que tenga la enfermedad, esté infectada y pueda transmitirla. Aquí entran los casos comunes de infección, pero también los casos asintomáticos, que en algunas ocasiones pueden incluso no tener registro.
    \item El compartimiento de recuperados incluye a las personas que sobrevivieron y salieron del compartimiento de infectados. Además, se supone que al recuperarse se adquiere inmunidad hacia la enfermedad, por lo que la probabilidad de que una persona recuperada pase a ser susceptible es nula.
\end{itemize}

En este modelo, la dinámica de transición entre los compartimientos, dependiente del tiempo, se define de la siguiente manera, con $N = S(t) + I(t) + R(t) + D(t)$:

\begin{align}
    \frac{dS}{dt} &= -\beta I \frac{S}{N} \\
    \frac{dI}{dt} &= \beta I \frac{S}{N} - \gamma I - \delta I \\
    \frac{dR}{dt} &= \gamma I \\
    \frac{dD}{dt} &= \delta I
\end{align}

Los parámetros que aparecen en las ecuaciones diferenciales son los encargados de modificar el comportamiento de la propagación de la enfermedad:

\begin{itemize}
    \item $\beta$ (beta): tasa de contagio. Controla qué tan rápido los susceptibles se infectan y depende del contacto social y de la probabilidad de transmisión.
    \item $\gamma$ (gamma): tasa de recuperación. $\frac{1}{\gamma}$ es el tiempo promedio de recuperación.
    \item $\delta$ (delta): tasa de mortalidad. $\frac{1}{\delta}$ es el tiempo promedio hasta el fallecimiento desde la infección.
\end{itemize}

Estos parámetros se toman como fijos, pero en la realidad no lo son: las distintas medidas o políticas públicas que se apliquen, así como el avance del conocimiento médico, afectan directamente su valor.

Ahora, conociendo lo que representa cada uno de los parámetros $\beta$, $\gamma$ y $\delta$, es fácil entender lo que describen las ecuaciones diferenciales. Sabemos que la derivada representa el cambio respecto a algo, que en este caso es el tiempo. Cada una de las ecuaciones nos dice cómo cambia el número de individuos en ese compartimiento con respecto al tiempo.

En el caso de los susceptibles, la cantidad de personas únicamente puede disminuir, razón por la cual aparece el signo menos. Además, estas personas salen de este compartimiento según la proporción de infectados y susceptibles respecto a la población total.

Por otro lado, las personas que salen del compartimiento de susceptibles son las mismas que entran al compartimiento de infectados, por eso el término se traslada directamente. Además, de este compartimiento hay dos vías de salida: la recuperación o la muerte.

Los dos últimos compartimientos son relaciones directas de lo anterior.


Estos parámetros, en los modelos clásicos, pueden tomarse como constantes; sin embargo, cuando se busca aplicarlos a un problema a gran escala y en tiempo real, dejan de serlo, ya que diversas acciones, como la aplicación de políticas públicas, el avance en el conocimiento de la enfermedad y el pánico mismo de la población, afectan directamente el valor de dichos parámetros.
Con este modelo, ahora podemos hacer estimaciones resolviendo numéricamente el sistema de ecuaciones. Esta resolución nos dirá cómo va avanzando la epidemia, dependiendo de los valores iniciales que nosotros coloquemos, además de los valores de los parámetros $\gamma$, $\beta$ y $\delta$.
Toda esta estimación será nuestra creencia.


%%%%%%%%%%%%%%%

Bajo estos supuestos, $S(t)+I(t)+R(t)+D(t)=N$ en todo instante. Cabe notar que la población viva es $N-D(t)$. Aquí se conserva $N$ en el denominador como referencia fija del modelo. Además, se supone una mezcla homogénea a escala de la población considerada, de modo que las diferencias de contacto entre individuos quedan resumidas en una sola tasa de contagio.
El primer término representa la cantidad de nuevas infecciones por unidad de tiempo. Es el mismo flujo que sale de $S$ y entra en $I$. Por otro lado, de $I$ salen dos flujos, uno hacia $R$ y otro hacia $D$. Así, cuando las tasas son no negativas, la cantidad de susceptibles solo puede disminuir o permanecer constante, mientras que las cantidades de recuperados y fallecidos son acumulativas. Si el tiempo se mide en días, las tres tasas se expresan en días inversos.

Para precisar la interpretación de los inversos de las tasas, es importante considerar que la recuperación y el fallecimiento son salidas que compiten entre sí. Cuando las dos tasas se mantienen constantes y su suma es positiva, el tiempo medio de permanencia en $I$ es $1/(\gamma+\delta)$ y la proporción que sale por fallecimiento es $\delta/(\gamma+\delta)$. Si las tasas cambian durante el periodo infeccioso, tanto la permanencia como el desenlace dependen de los valores que estas toman a lo largo de dicho periodo.

Con lo anterior podemos preguntarnos en qué condiciones crece el compartimiento de infectados. Para ello, definamos el número de reproducción instantáneo del modelo, siempre que $\gamma(t)+\delta(t)>0$, mediante
\begin{equation}
    \mathcal{R}_t=\frac{\beta(t)}{\gamma(t)+\delta(t)}\frac{S(t)}{N}.
\end{equation}
Cuando $I(t)>0$, su derivada es positiva si $\mathcal{R}_t>1$ y negativa si $\mathcal{R}_t<1$. La interpretación en términos de infecciones secundarias corresponde a mantener las condiciones de transmisión y remoción del instante considerado. Esta distinción es importante cuando dichas condiciones cambian durante el curso de una infección \cite{cori2013}\label{cita:cori2013}.

Ahora podemos preguntarnos qué representa un cambio en $\beta(t)$. Bajo el supuesto de mezcla homogénea, el riesgo instantáneo de infección de una persona susceptible está dado por $\beta(t)I(t)/N$. Por ello, cuando comparamos dos fechas, un riesgo menor puede deberse a que disminuyó la tasa de contagio, a que hay menos personas infecciosas o a una combinación de ambos efectos. La cantidad de nuevas infecciones depende, además, de cuántas personas sigan siendo susceptibles. Esta distinción permite entender por qué la curva de casos puede cambiar aun cuando algunas de las tasas permanezcan constantes.

De manera análoga, el flujo de defunciones depende de $\delta(t)$ y de la cantidad de personas en $I(t)$. Si se observa un aumento en las defunciones, una posible explicación es que haya aumentado el grupo infeccioso en fechas previas. Otra es que haya aumentado la tasa efectiva de salida hacia $D$. Para distinguirlas se necesita considerar la trayectoria de los demás compartimientos y la forma en que el modelo relaciona las infecciones con las defunciones. Tal relación, dentro del SIRD, se resume en la permanencia en $I$ y en las dos formas de salir de ese compartimiento.

Cabe notar que una tasa estimada para toda la población resume procesos que pueden ser diferentes entre grupos. Por ejemplo, al cambiar la composición del grupo infeccioso, puede variar la tasa agregada de mortalidad aun si el riesgo de cada subgrupo se mantiene. Por ello, cuando interpretemos las tasas como dependientes del tiempo, las entenderemos como cantidades efectivas del modelo para la población considerada. Atribuir una variación a una medida sanitaria o a un cambio en la composición de los casos requeriría información adicional.

Para comparar esta trayectoria con los registros, es necesario establecer primero qué mide cada serie. Por ejemplo, la incidencia que predice el modelo en un día corresponde a la disminución de $S$ entre los extremos de ese intervalo. De manera análoga, las defunciones que ocurren durante ese día corresponden al incremento de $D$. Los casos registrados por fecha de ingreso pueden reflejar infecciones ocurridas en días anteriores y dejar fuera a personas que nunca acudieron a una unidad de salud. Así, el modelo de observación debe representar, en la medida en que lo permitan los datos, la cobertura del registro y los desfases entre el evento y su observación.

Para precisar esta relación, consideremos un intervalo $(t_{k-1},t_k]$. Si la serie de defunciones se construye con la fecha de ocurrencia y el registro es completo, la cantidad comparable con el dato es $D(t_k)-D(t_{k-1})$. Si el dato es acumulado, la comparación puede formularse en términos de $D(t_k)$, siempre que el origen de la acumulación sea el mismo en ambas series. Hay que tener en cuenta que dos acumulados sucesivos comparten los eventos de las fechas anteriores. Por ello, cuando la incertidumbre proviene de errores en los conteos diarios que luego se acumulan, los errores de los acumulados sucesivos pueden estar correlacionados.

Por otro lado, para construir una serie de prevalencia se requiere saber quién permanece en el compartimiento infeccioso en cada fecha. Una fecha de ingreso permite ubicar un evento de registro, pero la duración de la estancia en $I$ depende de cuándo comienza y termina la capacidad de transmitir. Si solo se cuenta con fechas de ingreso, para aproximar esa prevalencia habría que introducir supuestos adicionales. Tales supuestos pasarían a formar parte del modelo de observación y su incertidumbre afectaría la estimación de las tasas.

Finalmente, que una trayectoria se ajuste a los registros no garantiza que las tasas se puedan determinar por separado. Por ejemplo, cuando $S(t)/N$ se aproxima a uno, el crecimiento de $I(t)$ informa sobre todo acerca de la diferencia $\beta(t)-\gamma(t)-\delta(t)$. Para distinguir los efectos de las tres tasas se requiere información adicional sobre los flujos o supuestos que restrinjan cómo pueden cambiar en el tiempo. Conviene distinguir la identificabilidad estructural, que depende del modelo y de las variables observadas en condiciones ideales, de la identificabilidad práctica, que también depende de la cantidad y la calidad de los datos disponibles \cite{roosa2019}\label{cita:roosa2019}.

Esta diferencia se aprecia si suponemos, como caso ideal, que pudiéramos observar sin error todos los compartimientos de forma continua. En los instantes en que $I(t)>0$, las pendientes de $R$ y $D$ permitirían determinar sus tasas de salida por separado. Si además $S(t)>0$, la pendiente de $S$ daría información directa sobre la tasa de contagio. Con registros diarios, en cambio, observamos cambios a lo largo de intervalos y solo de algunas variables. Estimar un valor instantáneo a partir de tales cambios requiere decidir qué variación de las tasas es razonable dentro de cada intervalo.

En particular, si durante varios días hay pocos individuos infecciosos, pueden observarse pocas recuperaciones o defunciones aun cuando las tasas sean positivas. Un ajuste que atribuya todos los cambios diarios a las tasas puede confundir la variación del registro con la dinámica de la epidemia. Por ello, cuando permitimos que las tasas dependan del tiempo, es necesario especificar también cuánto pueden cambiar de una fecha a otra. Esta especificación aporta información al problema y debe tomarse en cuenta al interpretar las estimaciones.

Esto nos lleva al problema que se desarrollará en el siguiente capítulo. Si solo observamos una parte de lo que ocurre en la epidemia, ¿cómo podemos actualizar el estado y las tasas a partir de esos datos? Para responder esta pregunta, plantearemos la evolución del sistema y la generación de las observaciones como dos partes de un modelo probabilístico.


%%%%%%%%%%%%%%%

